% TOP_PROBLEM: {"id": 20001034, "title": "A feedback-set bound for second neighborhoods", "category_id": 2, "category": {"id": 2, "name": "combinatorics", "display_name": "Combinatorics", "description": "Counting problems, graph theory, discrete structures.", "slug": "combinatorics", "order_index": 2, "created_at": "2026-07-31T15:26:25.671Z"}, "status": "partially_solved", "problem_number": "AIM-COMBINATORICS-0159", "set_id": 14}
% FIRST_POSTED: 2026-10-01
% ENTRY_KIND: statement_only
% UnsolvedMath ID 20001034; source catalogue code AIM-COMBINATORICS-0159
% !TeX program = lualatex
% Definition and sources only; this file is not an LLM solution attempt.
\documentclass[11pt,a4paper]{article}
\usepackage[margin=25mm]{geometry}
\usepackage{fontspec}
\setmainfont{lmroman10-regular.otf}[BoldFont=lmroman10-bold.otf,
 ItalicFont=lmroman10-italic.otf,BoldItalicFont=lmroman10-bolditalic.otf]
\usepackage{amsmath,amssymb,mathtools}
\usepackage{unicode-math}
\setmathfont{latinmodern-math.otf}
\AtBeginDocument{\let\setminus\smallsetminus}
\usepackage{xurl,hyperref}
\hypersetup{colorlinks=true,urlcolor=blue}
\setcounter{secnumdepth}{0}
\setlength{\parindent}{0pt}
\setlength{\parskip}{5pt}
\setlength{\emergencystretch}{3em}
\begin{document}
\section*{A feedback-set bound for second neighborhoods}\label{20001034}
{\small UnsolvedMath ID 20001034 · AIM-COMBINATORICS-0159 · Combinatorics · AIM Workshop Problem Lists}
\subsection{Definitions and mathematical statement}
An oriented graph is a finite simple directed graph with no loops and
no pair of opposite arcs. For a vertex $v$, its out-neighbourhood is
$N^+(v)=\{u:v\to u\}$, and $d^+(v)=|N^+(v)|$.
Its second out-neighbourhood is the set of vertices at directed distance
exactly two:
\[
 N_2^+(v)=\{w\notin N^+(v)\cup\{v\}:\text{there exists }u
                                   \text{ with }v\to u\to w\}.
\]
Distance means the length of a shortest directed path. In particular,
a vertex reached by both one and two arcs is excluded from $N_2^+(v)$,
and different two-edge paths to the same endpoint count only once.

Seymour's second neighbourhood conjecture asserts that every nonempty
oriented graph contains a vertex $v$ satisfying
\[
                            |N_2^+(v)|\ge d^+(v).
\]
The vertex may depend on the entire graph; the claim does not require
the inequality at every vertex or on average. No connectedness,
regularity, or tournament hypothesis is imposed. A sink, which has
outdegree zero, automatically satisfies the required inequality, so
only graphs without sinks can present an obstruction.
The condition excluding opposite arcs is essential to the stated
problem: a pair joined in both directions can have a nonempty first
out-neighbourhood and an empty second one. This formulation preserves
the exact-distance convention used in the AIM source throughout.
\subsection{Short English statement}
Must every finite oriented graph have a vertex with at least as many second out-neighbours as immediate out-neighbours?
\subsection{Sources}
\begin{itemize}
\item[S1] ULAM.AI and the UnsolvedMath Contributors, supplied problems.json, record 20001034 (AIM-COMBINATORICS-0159), AIM Workshop Problem Lists. Catalogue identity and source transcription; CC BY 4.0.
\url{https://huggingface.co/datasets/ulamai/UnsolvedMath}.
\item[S2] American Institute of Mathematics, The Caccetta-Haggkvist conjecture, item 3. Original workshop question underlying AIM-COMBINATORICS-0159.
\url{https://aimath.org/WWN/caccetta/caccetta.pdf}.
\end{itemize}
\end{document}
